% =============================================================================
% Block 3 --- Aula 3: Solow evidence (golden rule, technological progress, growth and TFP accounting).
% Weighted toward Prova 01 error cluster 2 (golden-rule algebra: carry the exponent, cross-check, knife edge).
\needspace{14\baselineskip}% keep the block heading, Q7's source line and its statement together
\bloco{3}{Solow continued: golden rule, TFP and growth accounting}{Aula 3 \,\textbullet\, Kurlat \S4.3--4.5, ch.~5 \,\textbullet\, 10 points}

\questao{(4 points) \textit{Source: Kurlat, P. (2020), \emph{A Course in Modern Macroeconomics}, Exercise 5.4 ``Measurement of the
Capital Stock'', pp.~94--95 --- statement reproduced verbatim, all items (a)--(d).}}

\noindent The capital stock is not easy to measure. A common way to try to do this is to:
\begin{itemize}[leftmargin=2.2em,topsep=2pt,itemsep=1pt]
\item guess a value for $K_0$ at some date in the past (maybe when statistics were first collected),
\item measure investment every year thereafter (which is also not that easy), and
\item use the equation
\[
K_{t+1}=(1-\delta)\,K_t+I_t \tag{5.5.1}
\]
to compute the capital stock for every year thereafter.
\end{itemize}
Consider an economy that is in a steady state without technological progress (and has been there for an infinitely long time).
Every year the level of investment is:
\[
I_t=0.2
\]
and the depreciation rate is:
\[
\delta=0.1
\]
\begin{enumerate}[label=(\alph*),leftmargin=2.2em,topsep=2pt,itemsep=1pt]
\item What is the capital stock?
\item The statistical office only started measuring investment in year 0. By then the economy was already in steady state. They
guessed (incorrectly) that the capital stock in year 0 was $K_0=1$. Use equation (5.5.1) to compute how the estimate of the capital
stock changed over time after year 0 and plot your answer. (You don't need to provide a closed form solution, a nice graph made
with a spreadsheet is enough). Compute the ratio:
\[
\frac{K_{ESTIMATE}}{K_{TRUE}}
\]
for the years 5, 10, and 50. (Again, you can just take the numbers from a spreadsheet).
\item Suppose now that the statistics office started collecting statistics a long time ago (an infinite time ago) so we don't have the
problem of making the wrong initial guess. However, instead of using the correct value of $\delta$, the statistical office incorrectly
believes that $\delta=0.05$. What is their estimate of the capital stock? Why does it differ from the truth?
\item The production function in this economy is
\[
Y=AK^{\alpha}L^{1-\alpha} \tag{5.5.2}
\]
and there are accurate data on GDP and on total hours of labor supply. Using the inaccurate estimate of $K$ from part (c), plus
accurate data on total labor and GDP, an economist is trying to measure the economy's productivity, i.e.\ to solve for $A$ in
equation (5.5.2). Find an expression for
\[
\frac{A_{ESTIMATE}}{A_{TRUE}}
\]
as a function of
\[
\frac{K_{ESTIMATE}}{K_{TRUE}}
\]
Given the numbers from part (c) and using $\alpha=0.35$, how far would the economist's estimate be from the truth?
\end{enumerate}
\noindent\textit{Changed by this list (not in the book):} for (b), give the closed form of the estimate, not only a spreadsheet.
\resolespaco
\begin{sol}
\begin{notacao}
\item $K_t$: capital stock; $I=0.2$: investment, constant every year; $\delta=0.1$: true depreciation rate.
\item $K^{TRUE}$: the true (steady-state) stock. $K^{EST}_t$: the statistician's estimate in year $t$, started from the guess $K^{EST}_0=1$.
\item $e_t\equiv K^{EST}_t-K^{TRUE}$: the estimation error. $\hat\delta=0.05$: the wrong depreciation rate of (c).
\item $A$: TFP, measured as a residual; $Y$, $L$: output and labour; $\alpha=0.35$: capital share.
\end{notacao}
\passo{(a) true stock}{Kurlat's perpetual-inventory equation (5.5.1), with constant investment:}
\begin{equation}
K_{t+1}=(1-\delta)K_t+I. \label{b3q7:pim}
\end{equation}
In a steady state $K_{t+1}=K_t=K$, so \eqref{b3q7:pim} gives
\begin{align*}
K&=(1-\delta)K+I\\
\delta K&=I\\
K^{TRUE}&=I/\delta=0.2/0.1=2 .
\end{align*}
\passo{(b) wrong initial guess}{The estimate and the truth both obey \eqref{b3q7:pim}. Subtracting one from the other, $I$ cancels:}
\begin{equation}
e_{t+1}=(1-\delta)\,e_t,\qquad e_0=1-2=-1. \label{b3q7:err}
\end{equation}
Iterating \eqref{b3q7:err} $t$ times:
\begin{align*}
e_t&=(1-\delta)^t e_0\\
&=(1-0.1)^t\,(-1)=-0.9^t,\\
K^{EST}_t&=K^{TRUE}+e_t\\
&=(2)+(-0.9^t)=2-0.9^t,\\
\frac{K^{EST}_t}{K^{TRUE}}&=\frac{(2-0.9^t)}{(2)}\\
&=1-\tfrac12\,0.9^t .
\end{align*}
Evaluating at the three dates:
\begin{center}\small
\begin{tabular}{@{}lccc@{}}
\toprule
Year $t$ & $0.9^t$ & $K^{EST}_t$ & $K^{EST}_t/K^{TRUE}$\\
\midrule
5 & 0.5905 & 1.4095 & 0.7048\\
10 & 0.3487 & 1.6513 & 0.8257\\
50 & 0.0052 & 1.9948 & 0.9974\\
\bottomrule
\end{tabular}
\end{center}
The guess is forgotten geometrically at rate $1-\delta$.
\passo{(c) wrong depreciation}{With $\hat\delta=0.05$ the perpetual inventory adds up every past vintage of investment, each
kept at the survival rate $1-\hat\delta$:}
\begin{align*}
K^{EST}&=\textstyle\sum_{j\ge0}(1-\hat\delta)^jI\\
&=I/\hat\delta=0.2/0.05=4,
\end{align*}
twice the truth: every vintage of investment is kept on the books twice as long as it really survives.
\passo{(d) TFP}{TFP is the residual of the production function:}
\begin{equation}
A=\frac{Y}{K^{\alpha}L^{1-\alpha}}. \label{b3q7:A}
\end{equation}
Dividing \eqref{b3q7:A} at $K^{EST}$ by \eqref{b3q7:A} at $K^{TRUE}$, with the same $Y$ and $L$:
\begin{align*}
\frac{A^{EST}}{A^{TRUE}}&=\frac{Y/\big[(K^{EST})^{\alpha}L^{1-\alpha}\big]}{Y/\big[(K^{TRUE})^{\alpha}L^{1-\alpha}\big]}\\
&=\Big(\frac{K^{EST}}{K^{TRUE}}\Big)^{-\alpha}\\
&=\Big(\frac{4}{2}\Big)^{-0.35}=2^{-0.35}=0.7846 :
\end{align*}
measured TFP is about $21.5\%$ too low.
\intuicao{TFP is a residual, so every error in measuring inputs lands in it with the opposite sign, damped by the input's share.}
\refbib{Kurlat (2020), \S5.4 (p.~86) and \S5.5 (p.~89). Consistent with the project's worked solution
(Problem 5.4, folder \texttt{Resolucao}); the book prints no answer.}
\rubrica{
\item (a) 0.5; (b) 1.5 (closed form 1, three ratios 0.5); (c) 1; (d) 1. Ratio in (d) written as $2^{\alpha}$ (sign of the bias reversed): max.\ 0.5 pt in (d).
}
\end{sol}

% -----------------------------------------------------------------------------
\questao{(5 points) Discursive. \emph{Golden rule $\times$ interest rates $\times$ dynamic efficiency.} Solow model in continuous time with
labour-augmenting technology: $Y=K^{\alpha}(AL)^{1-\alpha}$, $\dot A/A=g$, $\dot L/L=n$, depreciation $\delta$, saving rate $s$. Variables per
efficiency unit: $\tilde k=K/(AL)$, $\tilde y=\tilde k^{\alpha}$, $\tilde c=(1-s)\tilde y$, and $\dot{\tilde k}=s\tilde k^{\alpha}-(n+g+\delta)\tilde k$.
Competitive factor markets give $r=\alpha\tilde k^{\alpha-1}-\delta$. Illustrative calibration for Brazil: $\alpha=1/3$, $n=0.01$,
$g=0.02$, $\delta=0.05$, $s=0.18$.}

\textbf{a)} (2 points) Find steady-state $\tilde c^*(s)$ in closed form and the saving rate $s_{GR}$ that maximises it, with the second-order
condition. Cross-check with the capital-stock version of the problem, and use a knife-edge value of $\alpha$ to test your exponent.
\resolespaco
\begin{sol}
\begin{notacao}
\item $K$, $L$: capital and labour; $A$: labour-augmenting technology. $g\ge0$: growth of $A$; $n\ge0$: growth of $L$.
\item $\delta\in(0,1)$: depreciation. $s\in(0,1)$: saving rate. $\alpha\in(0,1)$: capital share.
\item $\tilde k\equiv K/(AL)$, $\tilde y$, $\tilde c$: capital, output and consumption per efficiency unit; a star marks the steady state.
\item $r$: real interest rate. $s_{GR}$: golden-rule saving rate. $\tilde k_{GR}$: golden-rule capital per efficiency unit.
\end{notacao}
\passo{the statement's equations}{Numbered here for reference:}
\begin{align}
\dot{\tilde k}&=s\tilde k^{\alpha}-(n+g+\delta)\tilde k, \label{b3q8:kdot}\\
\tilde y&=\tilde k^{\alpha},\qquad \tilde c=(1-s)\,\tilde y, \label{b3q8:yc}\\
r&=\alpha\tilde k^{\alpha-1}-\delta. \label{b3q8:r}
\end{align}
\passo{steady state}{Set $\dot{\tilde k}=0$ in \eqref{b3q8:kdot}:}
\begin{align}
s\tilde k^{\alpha}&=(n+g+\delta)\tilde k \label{b3q8:be}\\
\tilde k^{1-\alpha}&=\frac{s}{n+g+\delta}. \nonumber
\end{align}
Raise both sides to the power $1/(1-\alpha)$:
\begin{equation}
\tilde k^*=\Big(\frac{s}{n+g+\delta}\Big)^{\frac1{1-\alpha}}. \label{b3q8:kss}
\end{equation}
Substitute \eqref{b3q8:kss} into \eqref{b3q8:yc}:
\begin{align}
\tilde c^*(s)&=(1-s)(\tilde k^*)^{\alpha}\nonumber\\
&=(1-s)\Big(\Big(\frac{s}{n+g+\delta}\Big)^{\frac1{1-\alpha}}\Big)^{\alpha}\nonumber\\
&=(1-s)\Big(\frac{s}{n+g+\delta}\Big)^{\frac{\alpha}{1-\alpha}}. \label{b3q8:css}
\end{align}
The exponent is $\alpha/(1-\alpha)$, not $1/(1-\alpha)$: $\tilde k^*$ enters raised to $\alpha$.
\passo{first-order condition}{Take the log of \eqref{b3q8:css}:}
\begin{equation}
\ln\tilde c^*=\ln(1-s)+\frac{\alpha}{1-\alpha}\ln s-\frac{\alpha}{1-\alpha}\ln(n+g+\delta). \label{b3q8:ln}
\end{equation}
Differentiate \eqref{b3q8:ln} with respect to $s$ and set it to zero:
\begin{equation}
-\frac1{1-s}+\frac{\alpha}{(1-\alpha)s}=0. \label{b3q8:foc}
\end{equation}
Multiply \eqref{b3q8:foc} by $(1-\alpha)s(1-s)>0$:
\begin{align*}
\alpha(1-s)&=(1-\alpha)s\\
\alpha-\alpha s&=s-\alpha s\\
\alpha&=s .
\end{align*}
Hence
\begin{equation}
\boxed{s_{GR}=\alpha}. \label{b3q8:sgr}
\end{equation}
\passo{second-order condition}{Differentiate the left side of \eqref{b3q8:foc} once more:}
\[
\frac{d^2\ln\tilde c^*}{ds^2}=-\frac1{(1-s)^2}-\frac{\alpha}{(1-\alpha)s^2}<0 .
\]
At $s=\alpha$:
\begin{align*}
\frac{d^2\ln\tilde c^*}{ds^2}\Big|_{s=\alpha}&=-\frac1{(1-(\alpha))^2}-\frac{\alpha}{(1-\alpha)(\alpha)^2}\\
&=-\frac1{(1-\alpha)^2}-\frac{1}{(1-\alpha)\alpha}\\
&=-\frac{\alpha+(1-\alpha)}{\alpha(1-\alpha)^2}\\
&=-\frac1{\alpha(1-\alpha)^2}.
\end{align*}
A maximum.
\passo{cross-check}{Choose $\tilde k$ directly. By \eqref{b3q8:yc} and then \eqref{b3q8:be}, steady-state consumption is}
\begin{align}
\tilde c^*&=\tilde k^{\alpha}-s\tilde k^{\alpha}\nonumber\\
&=\tilde k^{\alpha}-(n+g+\delta)\tilde k. \label{b3q8:ck}
\end{align}
Maximising \eqref{b3q8:ck} over $\tilde k$:
\begin{equation}
\alpha\tilde k^{\alpha-1}=n+g+\delta, \label{b3q8:mpk}
\end{equation}
that is, MPK $=n+g+\delta$. Solve \eqref{b3q8:be} for $s$ and substitute \eqref{b3q8:mpk}:
\begin{align*}
s&=\frac{(n+g+\delta)\,\tilde k}{\tilde k^{\alpha}}\\
&=\frac{(\alpha\tilde k^{\alpha-1})\,\tilde k}{\tilde k^{\alpha}}\\
&=\alpha\tilde k^{\alpha-1}\,\tilde k^{1-\alpha}\\
&=\alpha .
\end{align*}
Same answer as \eqref{b3q8:sgr}.
\passo{knife edge}{As $\alpha\to0$ capital is useless, so $s_{GR}\to0$, as \eqref{b3q8:sgr} says. Dropping the $\alpha$ in the exponent
of \eqref{b3q8:css}, $\max(1-s)s^{1/(1-\alpha)}$ gives $s=1/(2-\alpha)\to1/2$: absurd, so the exponent is wrong.}
\refbib{Kurlat (2020), \S4.3 (p.~61) and \S4.5 (p.~69); Prova~01, Q2(b) (\texttt{Prova/}\allowbreak\texttt{prova-01-error-analysis.md}).}
\rubrica{
\item $\tilde c^*(s)$ with exponent $\alpha/(1-\alpha)$ 0.5; FOC and $s_{GR}$ 0.5; SOC 0.25; cross-check 0.5; knife edge 0.25.
Exponent $1/(1-\alpha)$ carried into the FOC: max.\ 0.75 pt. $s_{GR}$ confused with the $s$ that maximises $\tilde k^*$: max.\ 0.5 pt.
}
\end{sol}

\textbf{b)} (1.5 points) Compute $\tilde k^*$, $\tilde c^*$ and $r$ at $s=0.18$ and at $s_{GR}$. Show that $r-(n+g)$ has the sign of $\alpha-s$, and say
whether Brazil should raise its saving rate.
\resolespaco
\begin{sol}
\begin{notacao}
\item Symbols as in (a); calibration $\alpha=1/3$, $n=0.01$, $g=0.02$, $\delta=0.05$.
\end{notacao}
\passo{numbers}{Here $n+g+\delta=0.08$, $\frac{\alpha}{1-\alpha}=\frac12$ and $\frac1{1-\alpha}=\frac32$. Using \eqref{b3q8:kss} and \eqref{b3q8:css}
at $s=0.18$, where $s/0.08=2.25$:}
\begin{align*}
\tilde k^*&=\Big(\frac{0.18}{0.01+0.02+0.05}\Big)^{3/2}\\
&=2.25^{3/2}=3.375,\\
\tilde c^*&=(1-0.18)\Big(\frac{0.18}{0.01+0.02+0.05}\Big)^{1/2}\\
&=0.82\times2.25^{1/2}=0.82\times1.5=1.23 .
\end{align*}
At $s=1/3$, where $s/0.08=4.1667$:
\begin{align*}
\tilde k^*&=\Big(\frac{1/3}{0.08}\Big)^{3/2}\\
&=4.1667^{3/2}=8.5052,\\
\tilde c^*&=\Big(1-\frac13\Big)\Big(\frac{1/3}{0.08}\Big)^{1/2}\\
&=\tfrac23\times4.1667^{1/2}=\tfrac23\times2.0412=1.3608 .
\end{align*}
The ratio is $1.3608/1.23=1.106$: $10.6\%$ more consumption per efficiency unit.
\passo{interest rate}{Divide \eqref{b3q8:be} by $s\tilde k$:}
\begin{equation}
\tilde k^{\alpha-1}=\frac{n+g+\delta}{s}. \label{b3q8:kam1}
\end{equation}
Substitute \eqref{b3q8:kam1} into \eqref{b3q8:r}:
\begin{align}
r&=\alpha\Big(\frac{n+g+\delta}{s}\Big)-\delta\nonumber\\
&=\frac{\alpha(n+g+\delta)}{s}-\delta. \label{b3q8:rs}
\end{align}
Subtract $n+g$ from both sides of \eqref{b3q8:rs}:
\begin{align*}
r-(n+g)&=\frac{\alpha(n+g+\delta)}{s}-(n+g+\delta)\\
&=(n+g+\delta)\Big(\frac{\alpha}{s}-1\Big),
\end{align*}
positive iff $s<\alpha$. Using \eqref{b3q8:rs} at each saving rate:
\begin{align*}
s=0.18:\quad r&=\tfrac13\times\tfrac{0.08}{0.18}-0.05=0.1481-0.05=0.0981>0.03,\\
s=1/3:\quad r&=\tfrac13\times\tfrac{0.08}{(1/3)}-0.05\\
&=0.08-0.05=0.03=n+g .
\end{align*}
The golden rule is $r=n+g$.
\passo{policy}{Brazil is below the golden rule and dynamically efficient. Raising $s$ to $1/3$ raises long-run consumption, but
consumption falls on impact and for a while along the transition: generations are traded off, and Solow, with no preferences,
cannot say whether the sacrifice is worth it. That is the job of the microfounded models of Aulas 4--6.}
\refbib{Kurlat (2020), \S4.3 (p.~61) and \S4.4 (p.~63); Exercise~5.8 (p.~97).}
\rubrica{
\item Numbers 0.5; $r-(n+g)$ sign 0.5; policy reading 0.5. Concluding that being below the golden rule is inefficient: max.\ 0.25 pt for the reading.
}
\end{sol}

\textbf{c)} (1.5 points) Another economy has the same parameters but $s=0.45$. Show that cutting $s$ to $\alpha$ raises consumption per
efficiency unit on impact and at every later date, and name the inefficiency.
\resolespaco
\begin{sol}
\begin{notacao}
\item Symbols as in (a). $\tilde k_0$: capital per efficiency unit when $s$ is cut; $\tilde c_t$, $\tilde k_t$: values along the transition.
\item $f(\tilde k)\equiv\tilde k^{\alpha}-(n+g+\delta)\tilde k$: steady-state consumption as a function of $\tilde k$, the right side of \eqref{b3q8:ck}.
\end{notacao}
\passo{diagnosis}{At $s=0.45$, $s/0.08=5.625$. Using \eqref{b3q8:kss}, \eqref{b3q8:css} and \eqref{b3q8:rs}:}
\begin{align*}
\tilde k^*&=\Big(\frac{0.45}{0.08}\Big)^{3/2}\\
&=5.625^{3/2}=13.3409,\\
\tilde c^*&=(1-0.45)\Big(\frac{0.45}{0.08}\Big)^{1/2}\\
&=0.55\times5.625^{1/2}=0.55\times2.3717=1.3044<1.3608,\\
r&=\tfrac13\times\tfrac{0.08}{0.45}-0.05=0.0093<n+g=0.03 .
\end{align*}
\passo{impact}{$\tilde k$ is predetermined, so by \eqref{b3q8:yc} $\tilde y=2.3717$ on impact and $\tilde c$ jumps:}
\[
\tilde c:\ 0.55\times2.3717=1.3044\ \longrightarrow\ \tfrac23\times2.3717=1.5811 .
\]
\passo{transition}{Using \eqref{b3q8:yc} and then \eqref{b3q8:kdot} for $s\tilde k^{\alpha}$, consumption along the path is}
\begin{align}
\tilde c_t&=\tilde k_t^{\alpha}-s\tilde k_t^{\alpha}\nonumber\\
&=\tilde k_t^{\alpha}-\big(\dot{\tilde k}_t+(n+g+\delta)\tilde k_t\big)\nonumber\\
&=\tilde k_t^{\alpha}-(n+g+\delta)\tilde k_t-\dot{\tilde k}_t\nonumber\\
&=f(\tilde k_t)-\dot{\tilde k}_t. \label{b3q8:cpath}
\end{align}
With $s=\alpha$, $\tilde k$ falls monotonically from $13.3409$ to $\tilde k_{GR}=8.5052$, so $\dot{\tilde k}_t<0$. On
$(\tilde k_{GR},13.3409]$, $f$ is decreasing: $f'(\tilde k)=\alpha\tilde k^{\alpha-1}-(n+g+\delta)$ is negative above $\tilde k_{GR}$ by
\eqref{b3q8:mpk}. So $\tilde k_t<\tilde k_0$ gives $f(\tilde k_t)>f(\tilde k_0)$, and \eqref{b3q8:cpath} gives
\[
\tilde c_t>f(\tilde k_t)>f(\tilde k_0)=1.3044\quad\text{at every }t,
\]
ending at $1.3608$.
\passo{name}{Every date gains and none loses: a Pareto improvement. The economy over-accumulates capital, $r<n+g$:
\textbf{dynamic inefficiency}.}
\refbib{Kurlat (2020), \S4.3 (p.~61); Romer (2012), \S1.4 and \S2.11 (pp.~88--91).}
\rubrica{
\item Diagnosis 0.25; impact 0.5; transition argument 0.5; name 0.25. Asserting the path without the monotonicity argument: max.\ 0.25 pt for it.
}
\end{sol}

% -----------------------------------------------------------------------------
\questao{(1 point) Multiple choice. Illustrative Brazilian data, 2000--2014 averages: GDP grows $3.2\%$, the capital stock $4.4\%$ and hours
$1.8\%$ a year; the capital share is $0.4$. The Solow residual, and the share of the growth of output \emph{per hour} it accounts for, are:}
\begin{alternativas}
\item $0.36$ p.p.\ a year, accounting for about 11\% of the growth of output per hour.
\item $-0.16$ p.p.\ a year, accounting for about $-5\%$ of the growth of output per hour.
\item $1.44$ p.p.\ a year, accounting for about 45\% of the growth of output per hour.
\item $0.36$ p.p.\ a year, accounting for about 26\% of the growth of output per hour.
\item $0.56$ p.p.\ a year, accounting for about 40\% of the growth of output per hour.
\end{alternativas}
\gabarito{D}
\begin{sol}
\begin{notacao}
\item $g_Y$, $g_K$, $g_L$, $g_A$: growth of output, capital, hours and TFP (p.p.\ a year); $g_y$: growth of output per hour; $\alpha=0.4$: capital share.
\end{notacao}
\passo{residual}{Kurlat's growth-accounting equation (5.4.2):}
\begin{equation}
g_A=g_Y-\alpha g_K-(1-\alpha)g_L. \label{b3q9:ga}
\end{equation}
Substituting the data into \eqref{b3q9:ga}:
\begin{align*}
g_A&=3.2-0.4\times4.4-0.6\times1.8\\
&=3.2-1.76-1.08=0.36 .
\end{align*}
\passo{per hour}{Output per hour grows at $g_y=g_Y-g_L=3.2-1.8=1.4$. Subtracting $g_L$ from both sides of \eqref{b3q9:ga} and rearranging:}
\begin{align*}
g_y&=g_A+\alpha(g_K-g_L)\\
&=(0.36)+0.4\times(4.4-1.8)\\
&=0.36+0.4\times2.6=0.36+1.04=1.4 .
\end{align*}
The residual's share is $0.36/1.4=25.7\%$.
\passo{distractors}{A divides by output growth ($0.36/3.2$); B swaps the shares ($3.2-2.64-0.72$) and divides by $3.2$ (two mistakes);
C drops labour ($3.2-1.76$) and divides by $3.2$ (two); E swaps the shares and drops labour ($3.2-2.64$), divided by $1.4$ (two).}
\refbib{Kurlat (2020), \S5.4, eq.~(5.4.2) (p.~86).}
\end{sol}
